\documentclass[a4paper,11pt]{article}
\usepackage{exam-style-341}
\examinehead{341 农业综合知识三 · 网络技术分章练习}{第06章 网络层 参考答案与解析}
\begin{document}

\unitanswerscover{341 农业综合知识三 · 网络技术分章练习}{第06章 网络层 参考答案}{第06章练习的答案与解析；判断题均给出改正后的表述，计算题与综合题给出完整推导步骤，便于复算核对。}

\answerpart{一、单项选择题（10 题，每题 3 分，共 30 分）}

\q B。\begin{analysis}IPv4 数据报首部的前一部分长度固定，共 20 字节，是所有 IPv4 数据报必须具有的；首部长度字段最大可表示 60 字节（含可选字段）。A 8 字节是 UDP 首部长度；C 40 字节是 IPv6 基本首部长度；D 60 字节是 IPv4 首部含选项时的最大长度，不是固定部分。（教材 4.2.5，印刷页 136）\end{analysis}

\q C。\begin{analysis}片偏移指出某片数据在原数据报数据部分中的相对位置，以 8 字节为偏移单位，所以除最后一个分片外，各分片的数据长度必须是 8 字节的整数倍（教材 4.2.5，印刷页 137）。A 1 字节、B 4 字节、D 16 字节均不是片偏移的单位。\end{analysis}

\q B。\begin{analysis}私有地址有三段：10.0.0.0/8、172.16.0.0/12（即 172.16.0.0—172.31.255.255）、192.168.0.0/16。172.16.5.1 落在 172.16.0.0/12 内，是私有地址；A 172.15.5.1 不在 172.16—172.31 范围；C 172.32.5.1 也不在；D 192.169.5.1 不属于 192.168.0.0/16。\end{analysis}

\q C。\begin{analysis}C 类地址网络号占 3 个字节、主机号占 1 个字节，默认子网掩码为 255.255.255.0（/24）。A 255.0.0.0 是 A 类默认掩码；B 255.255.0.0 是 B 类默认掩码；D 255.255.255.255 是 /32 主机路由掩码。\end{analysis}

\q A。\begin{analysis}ARP 在同一局域网内通过广播 ARP 请求、单播 ARP 响应，把 IP 地址解析为 MAC 地址（教材 4.2.4，印刷页 133）。B ICMP 用于差错报告与询问；C IGMP 用于多播组管理；D BGP 是自治系统之间的路由协议。\end{analysis}

\q A。\begin{analysis}ICMP 报文作为 IP 数据报的数据部分封装在 IP 数据报中传送（协议字段值为 1），因此 ICMP 虽属于网络层协议，但并不直接封装在以太网帧或运输层报文中（教材 4.4）。B UDP、C TCP 属于运输层；D 以太网帧是链路层封装单位。\end{analysis}

\q B。\begin{analysis}RIP 以跳数为度量，规定一条路径最多只能包含 15 个路由器，跳数等于 16 时即认为目的网络不可达（教材 4.6.2，印刷页 159）。A 8、C 16、D 32 都不是 RIP 的最大跳数。\end{analysis}

\q B。\begin{analysis}OSPF 属于链路状态路由协议：各路由器洪泛发送链路状态通告，使全网建立一致的链路状态数据库，再用 Dijkstra 算法计算最短路径树（教材 4.6.3，印刷页 164）。A 距离向量是 RIP 的算法；C 路径向量是 BGP 的算法；D 随机路由不是标准路由算法。\end{analysis}

\q B。\begin{analysis}CIDR 斜线记法中，斜线后的数字表示网络前缀的位数，/26 表示前 26 位为网络前缀、后 6 位为主机号（教材 4.2.2，印刷页 126—127）。A 说反了；C 子网数不是由该数字直接给出；D 地址总数应为 $2^{32-26}=2^{6}=64$ 个，不是 26 个。\end{analysis}

\q C。\begin{analysis}IPv6 地址长度为 128 位，比 IPv4 的 32 位地址空间大得多（教材 4.5，印刷页 150）。A 32 位是 IPv4；B 64 位常见于 MAC 地址或 IPv6 接口标识；D 256 位不是 IPv6 地址长度。\end{analysis}

\answerpart{二、多项选择题（3 题，每题 4 分，共 12 分）}

\q ABC。\begin{analysis}IP 数据报首部固定部分包括版本、首部长度、区分服务、总长度、标识、标志、片偏移、生存时间、协议、首部检验和、源地址、目的地址等字段。D 错误：“目的端口”属于运输层（TCP／UDP）首部字段，不在 IP 首部中。\end{analysis}

\q ABC。\begin{analysis}私有地址范围为 10.0.0.0/8、172.16.0.0/12、192.168.0.0/16。A 10.1.2.3 在 10.0.0.0/8 内；B 172.16.0.1 在 172.16.0.0/12 内；C 192.168.1.1 在 192.168.0.0/16 内；D 172.32.0.1 已超出 172.31.255.255，不是私有地址。\end{analysis}

\q ABC。\begin{analysis}RIP、OSPF 是内部网关协议，BGP 是外部网关协议，三者都属于路由选择协议。D 错误：ARP 是地址解析协议，负责把 IP 地址解析为 MAC 地址，不参与路由选择。\end{analysis}

\answerpart{三、判断题（6 题，每题 2 分，共 12 分）}

\q ×。改正：IP 协议提供的是无连接的、尽最大努力交付的服务，不保证可靠交付、不保证按序到达，可靠性由上层协议（如 TCP）负责。\begin{analysis}IP 只负责把数据报从源主机送到目的主机，途中可能丢失、重复、乱序；把“面向连接、可靠”安到 IP 上是常见错误。\end{analysis}

\q √。\begin{analysis}划分子网是从主机号中借位作子网号，子网掩码中网络号和子网号对应的位为 1，因此 1 的个数比该类的默认掩码多。例如 C 类 202.117.35.0/24 借 3 位后掩码为 /27。\end{analysis}

\q √。\begin{analysis}路由聚合把若干连续的小地址块合并为一个较大的 CIDR 地址块，要求这些地址块连续且起始地址按聚合块大小对齐，否则聚合后的地址块会包含未分配的地址，造成路由错误。\end{analysis}

\q ×。改正：片偏移字段以 8 字节为单位，而不是以 1 字节为单位。\begin{analysis}片偏移占 13 位，表示分片数据在原数据报数据部分中的相对位置（以 8 字节为单位），因此非最后一片的数据长度应是 8 的整数倍。\end{analysis}

\q ×。改正：RIP 收敛速度慢，且最大跳数只有 15，只适合小型网络；大型网络通常采用 OSPF 等链路状态协议。\begin{analysis}RIP 采用距离向量算法，依靠相邻路由器周期性交换整张路由表，网络变化时“好消息传得快、坏消息传得慢”，收敛慢。\end{analysis}

\q √。\begin{analysis}默认路由 0.0.0.0/0 的斜线后数字为 0，表示网络前缀长度为 0，即掩码为 32 个 0，任何目的地址与它按位与的结果都是 0，故总能匹配；只有当其他表项都不匹配时才使用默认路由。\end{analysis}

\answerpart{四、计算题（1 题，共 6 分）}

\q 参考答案与推导：
（1）借位数与可用子网数（2 分）：可用子网数 $=2^{n}-2$，要求 $\ge 5$。$n=2$ 时 $2^{2}-2=2<5$，不够；$n=3$ 时 $2^{3}-2=6\ge 5$，故从主机号借 3 位作子网号，可用子网数为 6 个。
（2）子网掩码（2 分）：C 类默认前缀为 /24，借 3 位后前缀长度为 /27。掩码点分十进制为 255.255.255.224；二进制为
\[
11111111.11111111.11111111.11100000
\]
（3）每子网可用主机数（1 分）：主机号还剩 $8-3=5$ 位，可用主机数 $=2^{5}-2=30$ 台（去掉主机号全 0 的网络地址和全 1 的广播地址）。
（4）第 2 个可用子网（1 分）：子网块大小 $=2^{5}=32$。子网号（最后 3 位）依次取 001、010、011、100、101、110，对应网络地址最后一段为 32、64、96、128、160、192。第 1 个可用子网为 202.117.35.32/27，第 2 个可用子网为 202.117.35.64/27：
网络地址 202.117.35.64；广播地址 $=64+32-1=202.117.35.95$；可用主机地址范围 202.117.35.65—202.117.35.94（共 30 个）。
\begin{analysis}复核：$64=01000000_2$，最后 3 位 010 即第 2 个可用子网号；广播地址 95 的原因为主机位全 1（$64+31=95$）。若按现代设备支持全 0／全 1 子网计算，则可用子网数为 8 个，但本题按题干“经典划分子网方法”作答。\end{analysis}

\answerpart{五、简答题（2 题，每题 10 分，共 20 分）}

\q 参考答案要点（答出 8 个字段即可，每点 1 分，共 8 分；固定部分 20 字节、总述 2 分）：
（1）版本：占 4 位，表示 IP 协议版本（IPv4 为 4）。
（2）首部长度：占 4 位，以 32 位字（4 字节）为单位，最小 5（20 字节），最大 15（60 字节）。
（3）区分服务：占 8 位，用于获得更好的服务（旧称服务类型）。
（4）总长度：占 16 位，首部与数据之和，单位为字节，最大 65535 字节。
（5）标识：占 16 位，同一数据报的各分片具有相同的标识，用于重装。
（6）标志：占 3 位，含 DF（不能分片）和 MF（后面还有分片）两个有意义位。
（7）片偏移：占 13 位，以 8 字节为单位，指出该片数据在原数据报中的相对位置。
（8）生存时间 TTL：占 8 位，防止数据报在网络中无限制兜圈，每经一个路由器减 1。
（9）协议：占 8 位，指出数据部分使用何种上层协议（如 TCP=6、UDP=17、ICMP=1）。
（10）首部检验和：占 16 位，只检验首部，不检验数据部分，每经路由器重新计算。
（11）源地址、目的地址：各占 32 位，分别是发送和接收该数据报的主机 IP 地址。

\q 参考答案要点：
（1）RIP（3 分）：内部网关协议，采用距离向量算法。每个路由器维护到各目的网络的距离（跳数），周期性地把整张路由表发送给相邻路由器，按“最短距离”更新路由表；最大跳数 15，16 表示不可达。优点是简单，缺点是收敛慢、只适合小型网络。
（2）OSPF（3 分）：内部网关协议，采用链路状态算法。路由器洪泛发送链路状态通告，使全网建立一致的链路状态数据库，再用 Dijkstra 算法计算以本路由器为根的最短路径树；支持分层（区域）、收敛快，适合大型网络（教材 4.6.3）。
（3）BGP（2 分）：外部网关协议，采用路径向量算法，用于自治系统（AS）之间交换网络可达性信息，考虑策略而非单纯最短路径（教材 4.6.4）。
（4）比较（2 分）：RIP／OSPF 用于 AS 内部，BGP 用于 AS 之间；RIP 距离向量、收敛慢、规模小，OSPF 链路状态、收敛快、规模大，BGP 路径向量、面向策略与互联。

\answerpart{六、综合题（1 题，共 20 分）}

\q 参考答案与推导：
\textbf{（1）211.68.4.0/22 的地址块信息（6 分）}
前缀长度为 /22（1 分）；子网掩码为 255.255.252.0，即 11111111.11111111.11111100.00000000（2 分）；地址总数为 $2^{32-22}=2^{10}=1024$ 个（1 分）；地址范围为 211.68.4.0—211.68.7.255（2 分）。
理由：4 个连续 /24 地址块需要合并 $2^{2}=4$ 个，前缀长度为 $24-2=22$；第三字节 4—7 的二进制为 00000100—00000111，前 6 位相同（000001），后 2 位变化，故网络前缀到第三字节的前 6 位为止。

\textbf{（2）把 211.68.4.0/24 划分为 4 个等长子网（8 分）}
借位数：$2^{n}=4$，故借 $n=2$ 位（2 分）。子网掩码：/24 加 2 位得 /26，即 255.255.255.192（2 分）。主机位剩 $8-2=6$ 位，每个子网可用主机数 $=2^{6}-2=62$ 台（2 分）。第 1 个子网（子网号 00）为 211.68.4.0/26：网络地址 211.68.4.0，广播地址 $=0+63=211.68.4.63$，可用地址范围 211.68.4.1—211.68.4.62（2 分）。其余 3 个子网依次为 .64/26、.128/26、.192/26。

\textbf{（3）IP 数据报分片（6 分）}
每片数据部分最大长度 $=1500-20=1480$ 字节（1480 是 8 的倍数）（1 分）；$3000=1480+1480+40$，故分 3 片，数据长度分别为 1480、1480、40 字节（2 分）；片偏移分别为 $0/8=0$、$1480/8=185$、$2960/8=370$（2 分）；MF 标志分别为 1、1、0；总长度分别为 $1480+20=1500$、1500、$40+20=60$ 字节（1 分）。3 片的标识字段相同（复制自原数据报），目的主机依据“标识＋源地址＋目的地址＋协议”重装。
\begin{analysis}复核：第 3 片总长度为 60 字节而不是 40 字节（数据 40＋首部 20）；片偏移 370 是由数据起点 2960 除以 8 得到；MF 只有最后一片为 0。\end{analysis}

\end{document}
